% Options for packages loaded elsewhere
\PassOptionsToPackage{unicode}{hyperref}
\PassOptionsToPackage{hyphens}{url}
\documentclass[
  ignorenonframetext,
]{beamer}
\newif\ifbibliography
\usepackage{pgfpages}
\setbeamertemplate{caption}[numbered]
\setbeamertemplate{caption label separator}{: }
\setbeamercolor{caption name}{fg=normal text.fg}
\beamertemplatenavigationsymbolsempty
% remove section numbering
\setbeamertemplate{part page}{
  \centering
  \begin{beamercolorbox}[sep=16pt,center]{part title}
    \usebeamerfont{part title}\insertpart\par
  \end{beamercolorbox}
}
\setbeamertemplate{section page}{
  \centering
  \begin{beamercolorbox}[sep=12pt,center]{section title}
    \usebeamerfont{section title}\insertsection\par
  \end{beamercolorbox}
}
\setbeamertemplate{subsection page}{
  \centering
  \begin{beamercolorbox}[sep=8pt,center]{subsection title}
    \usebeamerfont{subsection title}\insertsubsection\par
  \end{beamercolorbox}
}
% Prevent slide breaks in the middle of a paragraph
\widowpenalties 1 10000
\raggedbottom
\AtBeginPart{
  \frame{\partpage}
}
\AtBeginSection{
  \ifbibliography
  \else
    \frame{\sectionpage}
  \fi
}
\AtBeginSubsection{
  \frame{\subsectionpage}
}
\usepackage{iftex}
\ifPDFTeX
  \usepackage[T1]{fontenc}
  \usepackage[utf8]{inputenc}
  \usepackage{textcomp} % provide euro and other symbols
\else % if luatex or xetex
  \usepackage{unicode-math} % this also loads fontspec
  \defaultfontfeatures{Scale=MatchLowercase}
  \defaultfontfeatures[\rmfamily]{Ligatures=TeX,Scale=1}
\fi
\usepackage{lmodern}
\ifPDFTeX\else
  % xetex/luatex font selection
\fi
% Use upquote if available, for straight quotes in verbatim environments
\IfFileExists{upquote.sty}{\usepackage{upquote}}{}
\IfFileExists{microtype.sty}{% use microtype if available
  \usepackage[]{microtype}
  \UseMicrotypeSet[protrusion]{basicmath} % disable protrusion for tt fonts
}{}
\makeatletter
\@ifundefined{KOMAClassName}{% if non-KOMA class
  \IfFileExists{parskip.sty}{%
    \usepackage{parskip}
  }{% else
    \setlength{\parindent}{0pt}
    \setlength{\parskip}{6pt plus 2pt minus 1pt}}
}{% if KOMA class
  \KOMAoptions{parskip=half}}
\makeatother
\usepackage{longtable,booktabs,array}
\usepackage{caption}
\captionsetup[table]{skip=6pt}
\newcounter{none} % for unnumbered tables
\usepackage{calc} % for calculating minipage widths
\usepackage{caption}
% Make caption package work with longtable
\makeatletter
\def\fnum@table{\tablename~\thetable}
\makeatother
\setlength{\emergencystretch}{3em} % prevent overfull lines
\providecommand{\tightlist}{%
  \setlength{\itemsep}{0pt}\setlength{\parskip}{0pt}}
% Manuscript macro/package fallbacks (newcommand -> providecommand).
\usepackage{amsmath}
\usepackage{amssymb}
\usepackage{booktabs}
% Contents listing: article.cls prefixes every \l@section entry with
% \addvspace{1.0em}, which across ten sections costs about ten lines and pushed
% the ~40-entry listing one line past page 2 — leaving a third sheet carrying a
% single "10 References" line. Tighten that inter-section gap for the duration
% of \tableofcontents only, inside a group, so body spacing is untouched and no
% entry is dropped (reducing \tocdepth would have hidden 29 real subsections).
  \providecommand{\addvspace}[1]{\vskip 2\p@}%

% Slide layout defaults for warning-clean scientific decks.
\usepackage{etoolbox}
\IfFileExists{xurl.sty}{\usepackage{xurl}}{}
\IfFileExists{seqsplit.sty}{\usepackage{seqsplit}}{\newcommand{\seqsplit}[1]{#1}}
\protected\def\breakseq#1{\seqsplit{#1}}
\protected\def\breaktt#1{\begingroup\ttfamily\seqsplit{#1}\endgroup}
\setlength{\emergencystretch}{6em}
\tolerance=5000
\hbadness=10000
\hfuzz=1pt
\setlength{\tabcolsep}{2pt}
\AtBeginEnvironment{longtable}{\tiny\renewcommand{\arraystretch}{0.86}\setlength{\tabcolsep}{1pt}}
\AtBeginEnvironment{tabular}{\tiny\renewcommand{\arraystretch}{0.86}\setlength{\tabcolsep}{1pt}}
\AtBeginEnvironment{equation}{\tiny}
\AtBeginEnvironment{equation*}{\tiny}
\AtBeginEnvironment{align}{\tiny}
\AtBeginEnvironment{align*}{\tiny}
\AtBeginEnvironment{itemize}{\footnotesize}
\AtBeginEnvironment{enumerate}{\footnotesize}
\AtBeginEnvironment{description}{\footnotesize}
\setbeamerfont{caption}{size=\tiny}
\setbeamerfont{caption name}{size=\tiny}
\setbeamerfont{normal text}{size=\small}
\setbeamerfont{frametitle}{size=\small}
\setbeamerfont{section title}{size=\footnotesize}
\setbeamerfont{subsection title}{size=\footnotesize}
\setbeamertemplate{section page}{%
  \centering
  \begin{beamercolorbox}[sep=12pt,center,wd=\paperwidth]{section title}
    \parbox{0.86\paperwidth}{\centering\usebeamerfont{section title}\insertsection\par}
  \end{beamercolorbox}
}
\setbeamertemplate{subsection page}{%
  \centering
  \begin{beamercolorbox}[sep=8pt,center,wd=\paperwidth]{subsection title}
    \parbox{0.86\paperwidth}{\centering\usebeamerfont{subsection title}\insertsubsection\par}
  \end{beamercolorbox}
}
\setlength{\abovecaptionskip}{2pt}
\setlength{\belowcaptionskip}{0pt}

% Natbib and cross-reference fallbacks — slides don't load natbib
% or cleveref, but combined-PDF manuscript prose may emit these
% commands. The fallback renders citations as a bracketed key list
% and cross-references as detokenized labels so slides don't fail on
% undefined control sequences or raw underscores. \providecommand is
% a no-op if packages load later.
\providecommand{\citep}[1]{[#1]}
\providecommand{\citet}[1]{#1}
\providecommand{\citealp}[1]{#1}
\providecommand{\citeauthor}[1]{#1}
\providecommand{\citeyear}[1]{#1}
\providecommand{\cref}[1]{\texttt{\detokenize{#1}}}
\providecommand{\Cref}[1]{\texttt{\detokenize{#1}}}
\providecommand{\crefrange}[2]{\texttt{\detokenize{#1}}--\texttt{\detokenize{#2}}}
\providecommand{\Crefrange}[2]{\texttt{\detokenize{#1}}--\texttt{\detokenize{#2}}}
\providecommand{\paragraph}[1]{\textbf{#1}\ }

\newtheorem{proposition}{Proposition}
\newtheorem{hypothesis}{Hypothesis}
\newtheorem{remark}[theorem]{Remark}
\newtheorem{axiom}{Axiom}
\newtheorem{property}{Property}
\usepackage{bookmark}
\IfFileExists{xurl.sty}{\usepackage{xurl}}{} % add URL line breaks if available
\urlstyle{same}
% fallback for those not using the hyperref driver hyperxmp:
\makeatletter
\@ifundefined{xmpquote}{\newcommand{\xmpquote}[1]{#1}}{}
\makeatother
\hypersetup{
  hidelinks,
  pdfcreator={LaTeX via pandoc}}

\author{\texorpdfstring{}{}}
\date{}

\begin{document}

\section{Introduction}\label{sec:introduction}

\begin{frame}[allowframebreaks]{Motivation}
\protect\phantomsection\label{motivation}
Active Inference has matured into a broad research program spanning the
free energy principle (Friston 2010), discrete state-space formulations
of perception and action (Da Costa et al. 2020), and a growing body of
pedagogy and reference implementations (Parr et al. 2022; Smith et al.
2022). Recent graphical-specification and message-passing work makes the
same practical pressure visible: active-inference agents need
representations that can be inspected as model structure and then
carried into executable inference updates without being reauthored by
hand (Koudahl et al. 2023; Laar et al. 2023; Bagaev and Vries 2021). As
the field has grown, so has a quieter problem: the models themselves are
difficult to share, reproduce, and compare. A generative model that
lives only as a tangle of matrix definitions inside one author's script,
a diagram in a slide deck, and a paragraph of prose in a paper has no
\end{frame}

\begin{frame}[allowframebreaks]{Motivation}
single authoritative form. The same model is described three times, in
three incompatible media, with no guarantee that they agree. When a
reader wants to re-run that model in a different toolkit, they must
reverse-engineer it from whichever fragment they happen to have.

This is a reproducibility and interoperability crisis specific to the
structure of Active Inference work. The discipline depends on precisely
specified state spaces, observation and transition tensors, prior
preferences, and policy structures; small notational ambiguities
propagate into materially different behaviour. Yet the community has
lacked a notation that is simultaneously human-readable,
machine-parseable, and faithful to the underlying mathematics.
Generalized Notation Notation (GNN) was introduced to fill exactly that
gap (Smékal and Friedman 2023): a standard, text-based way to write down
an Active Inference generative model once, such that every downstream
\end{frame}

\begin{frame}[allowframebreaks]{Motivation}
use derives from the same source of truth.
\end{frame}

\begin{frame}[allowframebreaks]{The GNN Approach}
\protect\phantomsection\label{the-gnn-approach}
GNN treats the model specification as a first-class artifact. A model is
written in a plain-text language whose syntax captures the components an
Active Inference modeller actually reasons about --- state factors,
observation modalities, the matrices that link them, control structure,
and the temporal organization of the model. Because the specification is
text, it lives comfortably in version control, diffs cleanly across
revisions, and can be authored and reviewed by humans without
specialized tooling.

\end{frame}

\begin{frame}[allowframebreaks]{The GNN Approach}
The defining commitment of GNN is what the project calls the Triple
Play: a single text specification is the common origin for three
coordinated renderings of the same model. The text form is the
authoritative, editable description. From it, GNN produces graphical
visualizations that expose the model's factor and dependency structure
for inspection and communication. And from the same source it produces
executable cognitive models that can actually be run, so that the
diagram, the equations, and the running code are provably the same model
rather than three artifacts that merely claim to be. The Triple Play
turns a model from a description scattered across media into one
specification with multiple faithful projections (see
fig.~2).
\end{frame}

\begin{frame}[allowframebreaks]{Contributions}
\protect\phantomsection\label{contributions}
This work contributes a standard notation together with the
infrastructure that makes it usable in practice:

\end{frame}

\begin{frame}[allowframebreaks]{Contributions}
\begin{itemize}
\tightlist
\item
  \textbf{A parseable, human-readable notation} for Active Inference
  generative models, whose text form is authoritative and from which
  every other representation is derived.
\item
  \textbf{A 25-step processing pipeline}, implemented across 44 source
  packages (0--24), that carries a specification from parsing and
  validation through visualization, rendering, and execution.
\item
  \textbf{9 rendering backends} (PyMDP, RxInfer.jl, ActiveInference.jl,
  JAX, DisCoPy, PyTorch, NumPyro, Stan, bnlearn) that materialize a
  single GNN specification as executable cognitive models across
  multiple simulation frameworks, realizing the executable arm of the
  Triple Play.
\item
  \textbf{141 Model Context Protocol tools} that expose the pipeline's
  capabilities to agentic and programmatic clients, so the notation and
  its tooling are directly accessible to automated workflows.
\item
  \textbf{9-family reliability gates} that exercise the pipeline against
  a curated set of model families (basics, discrete, continuous,
  hierarchical, multiagent, precision, structured, gridworld,
  scaling-study), turning interoperability claims into checks that must
  pass rather than assertions that are merely made.
\end{itemize}
\end{frame}

\begin{frame}[allowframebreaks]{Reader Orientation}
\protect\phantomsection\label{reader-orientation}
The remainder of this manuscript is organized to move from language to
mechanism to evidence. sec.~3 states the notation
and the generative model it denotes, writing the factorization, the four
matrices, and the two free-energy objectives as explicit equations
(eq.~1 through eq.~8); it presents
the pipeline structure (fig.~1), the Triple Play
(fig.~2), and the per-step responsibility table
(tbl.~1). sec.~4 follows a
specification through the pipeline --- parsing, type checking,
rendering, execution, and reporting --- and presents the backend
registry (tbl.~2), the backend capability
surface (fig.~3), and the long-running
orchestration contracts (fig.~4).
sec.~5 reports what the project has produced
and how each number is grounded, with the model-family registry
(tbl.~3), the family-by-framework coverage
(fig.~5), and repository-scale metrics
\end{frame}

\begin{frame}[allowframebreaks]{Reader Orientation}
(fig.~6). Claims of independent re-execution are
addressed in sec.~6, the boundaries of the
current system in sec.~7, the symbol and
construct tables in sec.~9, and full source
details in sec.~10. Read in this order, the manuscript
moves from why a standard notation is needed, through how GNN realizes
it, to the evidence that the standard holds.
\end{frame}

\end{document}
